\subsection{Checks for \texorpdfstring{Section~\ref{sec:sca}}{scale}}\label{chk:sca}
% src: saved outputs of the L4, B3, B5, L7 and C2 computations, as described in each check

These computations confirm results of \cref{sec:sca} on small instances. Where a second, independent computation was made, both are reported.

\begin{numcheck}[A moment bound at spectral radius one]\label{chk:sca-relations}
For the lazy walk of \cref{lem:sca-relations}(c) on $\Z^3$, quadrature of $G(E)=\int_0^\infty e^{-Es}I_0(s)^3\mathrm ds$ gave $G(3)=0.505462$ and $G(5)=0.214294$, hence the variance $2(G(3)+G(5))q=1.439512q$. The mean of $(1-\hat W(p)^2)^{-1}$ over midpoint Fourier grids with 40 and 80 points per axis was $1.4256$ and $1.4326$, and the partial sums of $\sum_m(2\pi)^{-3}\int\hat W^{2m}\mathrm dp$ on a grid with 48 points per axis were $1.329$, $1.404$ and $1.427$ when summed through $m=10$, $100$ and $1000$, increasing and bounded. For the lazy nearest-neighbor walks on $\Z$ ($W_{xx}=1/2$, neighbor weights $1/4$) and on $\Z^2$ ($W_{xx}=1/3$, neighbor weights $1/6$) the grid means grew without bound as the grid was refined, $25.4$, $100.4$ and $400.4$ on $\Z$ and $2.75$, $3.41$ and $4.07$ on $\Z^2$ at 50, 200 and 800 points per axis.
\end{numcheck}

\begin{numcheck}[Light cone of a nonlinear rule]\label{chk:sca-lightcone}
The synchronous update $\mu_i\leftarrow\tanh\bigl(\sum_kB_{ik}\mu_k\bigr)+h_i$ on a ring of 30 agents, with $B_{ii}$ uniform on $[0.2,0.6]$, the two neighbor entries of each row uniform on $[-0.5,0.5]$, offsets $h_i$ normal with standard deviation $0.1$, a standard normal initial state, and the influence bound $\bar J_{ik}=|B_{ik}|$. The Jacobian $J_{ij}(t,0)$ vanished outside the cone $d(i,j)\le t$ to machine precision at $t=3$, $6$ and $12$, and no entry exceeded $(\bar J^t)_{ij}$. The largest row sums of $\bar J^t$ were $1.51$, $1.93$ and $3.03$, against $g_{\mathrm{tot}}^t=2.10$, $4.39$ and $19.3$.
\end{numcheck}
% src: L4_lr Part A (numpy default_rng seed 12345; draws in the order B row by row, h, initial state)

\begin{numcheck}[Poisson clocks]\label{chk:sca-poisson}
A ring of 40 agents with clock rate one and linear scalar updates with coefficient matrix $A$. For random convex weights with largest diagonal entry $\bar J_{\mathrm{diag}}=0.726$ and largest off-diagonal row sum $g=0.762$, the exact mean $e^{\nu t(A-I)}$ never exceeded the bound of \cref{prop:sca-poisson} at $t=2$, $5$ and $10$, the largest ratio falling from $0.887$ to $0.453$; with constant diagonal $0.5$ the ratio was $1.000000$ at all three times, the predicted equality. For signed weights a dominating process with weights $|A|$ was never exceeded pathwise, and the Monte Carlo excess over the bound stayed within one standard error. For the directed copy chain the expected influence $\E J_{ij}(t)=e^{-\nu t}(\nu t)^d/d!$ equals the middle member of \eqref{eq:sca-poisson}, $0.17547$ at $t=5$, $d=5$ and $0.12511$ at $t=10$, $d=10$, from the matrix exponential. The probability that a signal held at the source has crossed $d$ hops by time $t$ is a different quantity and equals the last member $e^{-\nu t}\sum_{n\ge d}(\nu t)^n/n!$, $0.55951$ and $0.54207$ at the same points, with Monte Carlo values $0.555$ and $0.542$. The velocity $\inf_{\varphi>0}(ge^\varphi+\bar J_{\mathrm{diag}}-1)/\varphi$ evaluated to $2.7183$ at $g=\bar J_{\mathrm{diag}}=1$ and to $1.0001$ at $g=1$, $\bar J_{\mathrm{diag}}=0$.
\end{numcheck}
% src: L4_lr Part B

\begin{numcheck}[Belief-dependent routing and the switch chain]\label{chk:sca-routing}
Three hundred random routing instances with five policies, linear updates and linear expected free energy: the largest ratio of the exact Jacobian of the averaged map to $\bar J(\gamma)$ was $0.572$, with no violation. In the chain of \cref{ex:sca-switch-chain}, the response 20 hops away to a perturbation of $10^{-6}$ at the head equaled $(\gamma\alpha_0/2)^{20}\cdot10^{-6}$ for $\gamma\alpha_0=1$, $1.8$, $2$ and $2.5$ (for example $9.537\times10^{-13}$ at $\gamma\alpha_0=1$), was $0.734$ for $\gamma\alpha_0=4$, where the linear value $1.049$ lies outside the range of the map and the response approaches the fixed point $\mu^\ast=0.9575$ over further hops ($0.9575$ at 25 hops), and was $1.0$ at $\gamma=\infty$. With sampled routing and message noise of standard deviation $0.5$, quadrature gave the flip probabilities $0.15546$, $0.06767$, $0.03430$, $0.02345$ and $\Pr(Z<-2)=0.02275$, and the correlation lengths of the routing bit $2.69$, $6.88$, $14.07$, $20.82$ and $21.47$ hops, for $\gamma\alpha_0=2$, $4$, $8$, $32$ and $\infty$; Monte Carlo estimates from $4\times10^5$ draws agreed within two standard errors. Simulated correlations of the routing bit agreed with $(1-2p)^n$, for example $0.478$ against $0.483$ at $n=5$ and $\gamma\alpha_0=4$.
\end{numcheck}
% src: L4_lr Part C; flip probabilities recomputed by quadrature in this editing round (the Monte Carlo values of L4_lr were 0.1558, 0.0678, 0.0344, 0.0236, 0.0231)

\begin{numcheck}[One common code]\label{chk:sca-common-code}
This check collects the evidence for \cref{conj:sca-common-code}. The instances were drawn as follows. The two-agent loops live in a two-dimensional world with $F=fI$, $Q=I$ and unit message and sensing noise variances, with $f=0.5+0.49/(1+e^{-x})$ for $x$ normal with standard deviation $2$, and with code and sensing rows of independent normal entries with standard deviation $3$. The three-dimensional instances have $f$ uniform on $[0.5,0.97]$, $Q=I$ and code and sensing matrices of independent normal entries with standard deviation $2$, 120 draws for each of the rank pairs $(1,1)$, $(2,1)$ and $(1,2)$. The switching instances have $f$ uniform on $[0.8,0.995]$, $Q=qI$ with $q$ uniform on $[0.05,1]$, and rank-one code and sensing rows of normal entries with standard deviations $3$ and $6$. In random two-agent reciprocal loops that sense and listen, the largest ratio $\rho(W)/f$ was $1.0000000000$ over 300 draws and after Nelder--Mead refinement with a common code, against 100 of 300 draws above one with independent codes (largest $1.1673$, and $1.2071$ after refinement). In three-dimensional worlds with code and sensing ranks one and two, none of 360 common-code instances exceeded one, against 85 of 360 with independent codes (largest $1.1772$). For period-two switching between listening and sensing, none of 300 common-code instances exceeded $f$ per round, against 112 of 300 with independent codes (largest $1.0941$). The common-code instances with symmetric listening that confirm \cref{prop:sca-sufficient}(b) and the comparison of \cref{cor:sca-size}(b), and the one-round reception maps on random directed listening graphs, none of which exceeded spectral radius one, are those of \cref{chk:min-common-code}; they are reported only there.
\end{numcheck}
% src: L7-C1 evidence (L7_cexb_codes Parts B, C, E, G; numpy default_rng seed 20261108, parts run in the order of the script); B3_hyp and L7_cexb_codes Part F are reported in chk:min-common-code

\begin{numcheck}[Clustering without contraction]\label{chk:sca-no-moment}
Rings of 40 scalar agents over 18 rounds from a diagonal initial covariance, in four cases: a time-invariant signed $W$ with $g_0=1.52$ and spectral radius $0.77$; a time-varying signed schedule with time-varying noise, $g_0=1.49$; a nonnegative $W$ with every row sum above one, $g_0=1.60$ and spectral radius $1.18$; and a signed $W$ with spectral radius $1.33$, $g_0=2.22$. In each case the self weight and the two neighbor weights of an agent have moduli uniform on $[0.2,1]$ times a common scale ($0.55$, $0.5$, $0.62$ and $0.8$ in the four cases), with signs flipped independently with probability one half in the signed cases; the noise is diagonal with variances uniform on $[0.5,1.5]$, redrawn every round in the time-varying case, and the initial variances are uniform on $[0,2]$. Every covariance entry with $d(i,j)>2t$ was exactly zero. Over all pairs and rounds, the largest ratio of $\Vert\Sigma_{ij}(t)\Vert$ to the bound of \cref{thm:sca-no-moment}(a) was $0.42$ at $\varphi=0$, $0.23$ at $\varphi=0.5$, $0.20$ at $\varphi=1$ and $0.23$ at $\varphi=2$, and the normalized ratio of (b) was at most $0.06$. An independent computation on a ring of 40 agents with a signed time-varying schedule ($g_0=2.67$; self and nearest-neighbor entries uniform on $[-0.9,0.9]$ and noise variances uniform on $[0.2,1]$, both redrawn every round, initial variances uniform on $[0.5,1.5]$) over 16 rounds found no nonzero entry beyond $2t$ and largest ratios $0.207$, $0.106$ and $0.105$ at $\varphi=0$, $0.5$ and $1$. The lag bound \eqref{eq:sca-no-moment-lag} was tested separately on rings of 24 scalar agents over 12 rounds with lags $h=1,\dots,4$, for a time-invariant signed $W$ ($g_0=2.41$, spectral radius $1.42$), a time-varying signed schedule with time-varying noise ($g_0=3.79$) and a nonnegative $W$ with every row sum above one ($g_0=1.63$, spectral radius $1.43$); the self and nearest-neighbor entries are normal with standard deviation $0.5$ in the signed cases and uniform on $[0.3,0.6]$ in the nonnegative case, the initial variances are uniform on $[0.2,1]$, and the noise covariance is the identity except in the time-varying case, where its diagonal is uniform on $[0.3,1]$ and redrawn every round. Every entry of $\Cov(e(t+h),e(t))$ with $d(i,j)>2t+h$ was exactly zero, entries with $2t<d(i,j)\le2t+h$ were nonzero, and the largest ratio of an entry to the bound was $0.25$ at $\varphi=0$, $0.15$ at $\varphi=0.5$ and $0.08$ at $\varphi=1$. On the scalar chain with $W_{xx}=0.2$, $W_{x,x\pm1}=0.3$, $q=1$ and zero initial covariance, the lagged covariance at $t=2$, $h=1$, $d=3$ is $0.027$, against the bound $0.512$ at $\varphi=0$.
\end{numcheck}
% src: B3_nomoment Part A (numpy default_rng seed 20261108); C2_b3_nomoment (seed 17); lag test fix/08/scripts/sca1-01.py (seed 801)

\begin{numcheck}[Rank-one dominance in the amplifying profile]\label{chk:sca-rank-one}
The profile of \cref{ex:sca-amplifying-profile} rebuilt from its specification has dominant eigenvalue $1.0613266750560013$, simple and real, second modulus $0.412032$, and block row sums $1.0999$ for both agents. From the covariance recursion of its mean map, the per-round growth $\frac12\ln(\Vert\Sigma(t)\Vert/\Vert\Sigma(t-1)\Vert)$ was $0.06294$, $0.05969$, $0.05952013$ and $0.05951971$ at $t=25$, $50$, $100$ and $150$, against $\ln\rho=0.05951971$. One minus the top canonical correlation of the two agents was $1.6\times10^{-2}$, $8.0\times10^{-4}$, $2.1\times10^{-6}$, $5.4\times10^{-9}$ and $1.4\times10^{-11}$ at $t=25$, $50$, $100$, $150$ and $200$, evaluated in 60-digit arithmetic because double precision loses the last value to cancellation; the last three values fall by the factor $\rho^{-100}=2.6007\times10^{-3}$ per 50 rounds, the rate $\rho^{-2t}$ of \eqref{eq:sca-rank-one}.
\end{numcheck}
% src: B3_nomoment Part B

\begin{numcheck}[An amplifying site]\label{chk:sca-amplifying-site}
An open chain of 801 agents with $c=a=0.3$ ($\lambda=0.9$) and an amplifying center. For $\rho=1.1$ and $1.2$ the symmetric eigensolver reproduced $\rho$ to ten digits, the second eigenvalue was $0.899982\le0.9$, and the fitted eigenvector decay was $0.79536555$ against $m_\rho=0.79536546$ and $0.96242368$ against $0.96242365$. From the direct recursion $\Sigma(t+1)=W\Sigma(t)W^\top+qI$, the fronts at $\theta=0.5$ and $0.9$, resolved to whole hops, stood at $8$, $19$, $29$ and $6$, $17$, $28$ at $t=60$, $150$, $240$ for $\rho=1.1$, with fitted speeds over $t\ge150$ of $0.1152$ and $0.1182$ against $v_{\mathrm f}=0.1198$; for $\rho=1.2$ the speeds were $0.1946$ and $0.1855$ against $0.1894$. The lower bound of \cref{prop:sca-amplifying-site}(c) held at all 340 and 382 pairs $(t,x)$ checked with $\mathfrak s_x>1$, and the decomposition of (b) agreed with the direct recursion at $t=40$ to $2.5\times10^{-14}$. An independent computation on a chain of 401 agents found the top eigenvalue equal to $c+(4a^2+b_0^2)^{1/2}$ to ten digits, the second eigenvalue $0.899927$, the eigenvector decay $0.75848615$ against $\arcsinh(b_0/(2a))=0.75848614$, $\Vert B(t)P_v\Vert=1.3\times10^{-11}$, and the spectrum of $B(t)$ on $v^\perp$ in $[1.0000,5.2584]\subset[q,q/(1-\lambda^2)]=[1,5.2632]$.
\end{numcheck}
% src: B3_nomoment Part C; C2_b3_nomoment

\begin{numcheck}[Spectral form with signed weights]\label{chk:sca-spectral}
Three symmetric rings of 80 agents with signed weights, spectral radius $0.95$ and largest row sums $1.227$, $1.276$ and $1.271$, so that the row-sum bound is unusable. Each ring has diagonal entries uniform on $[-0.3,0.3]$ and symmetric bond weights uniform on $[-0.5,0.5]$, rescaled to spectral radius $0.95$, with noise covariance $qI$. The result: no violation of the bound of \cref{prop:sca-spectral}, with largest ratios $0.508$, $0.478$ and $0.539$.
\end{numcheck}
% src: L4_oz (d) (numpy default_rng seed 7, created at part (d) and continued into part (e))

\begin{numcheck}[Diffusive clustering]\label{chk:sca-diffusive}
A $20\times20$ torus with symmetric bond conductances uniform on $[0.5,1.5]$ and self-conductances uniform on $[0.5,2]$, with $P$ the conductances divided by their row sums, $W=(1-\varepsilon)P$ and noise covariance $qI$, at $\varepsilon=0.05$ and $0.01$: no violation of \eqref{eq:sca-diffusive}, with the bound minimized over $\vartheta\in\{0.05,0.10,\dots,0.95\}$. The fitted decay rates of the largest covariance over hop distances 4 to 17 were $0.390$ and $0.153$, against the guaranteed rates $0.320$ and $0.142$ and the row-sum rates $0.051$ and $0.010$. On a $15\times15$ torus with bond conductances uniform on $[0.2,1]$, self-conductances uniform on $[0.5,2]$, $\lambda=0.97$ and $q=1$, $|\Sigma_{ij}|$ used at most $2.3$ percent of the bound.
\end{numcheck}
% src: L4_oz (e) (generator of part (d), seed 7, continued); C_v09 (seed 23)

\begin{numcheck}[Heterogeneous media]\label{chk:sca-heterogeneous}
Rings of 300 agents with $F=0.97$ and nearest-neighbor attention, the ring form of the upper bound being computed by image sums. Homogeneous control: $m_{\mathrm{up}}=m_{\mathrm{low}}=$ fitted rate $=0.48879$. Mild heterogeneity (leaks in $[0.0495,0.0591]$, couplings in $[0.233,0.252]$): the guaranteed rate, computed here and in the cases below as the largest tilt $\varphi$ with $\Psi(\varphi)<1$ in \cref{prop:sca-heterogeneous}(a), was $0.4424$, slightly above the closed form $m_{\mathrm{up}}=\arccosh(1+\varepsilon_{\min}/(2a_{\max}))=0.4399$ and below the fitted rates $0.4667$ and $0.4699$ of the largest and smallest correlation at each distance, which are below $m_{\mathrm{low}}=0.6296$, so $\xi\in[1.59,2.26]$ hops. Strong symmetric heterogeneity: largest tilt $0.3591$ (closed form $0.3405$) $\le0.4795,0.5057\le1.5920$. Strong directed heterogeneity, where the two neighbor weights of an agent differ and the closed form does not apply: largest tilt $0.1599\le0.4906,0.5649\le2.1048$. The upper bound was used at most to $0.21$ on all resolved entries, $\Sigma-q_{\min}(I-W_{\mathrm{low}}^2)^{-1}\ge-1.9\times10^{-15}$ (round-off), and $\min_i(1-\varepsilon_i)\le\rho(W)\le\max_i(1-\varepsilon_i)$ in every case. On an open chain the rows of the two one-sided end agents lower the largest admissible tilt to $0.2644$, against $0.3591$ on the ring drawn from the same law, and the bound with $\bar\psi$ taken over all rows holds with largest ratio $0.11$; on an open $18\times18$ grid both bounds hold along the axis and the diagonal (largest ratio $0.045$, smallest gap $4.1\times10^{-7}$); on a ring with a period-three schedule of heterogeneous directed gains, after 600 rounds the largest ratio was $0.071$.
\end{numcheck}
% src: B3_hetero

\begin{numcheck}[The ring formula]\label{chk:sca-oz-ring}
The closed form \eqref{eq:sca-oz-1d} against a Lyapunov solution: relative difference at most $2.1\times10^{-15}$ at $(c,\varepsilon)=(0.25,0.2)$, $(0.25,0.02)$, $(0.6,0.01)$ and $(0.05,0.05)$, and, with ring images, absolute difference at most $1.9\times10^{-13}$ on rings of 50 to 100 agents at four further parameter sets, one with $c=0$. At $(c,\varepsilon)=(0.25,0.02)$ the rate is $m=0.23355=2\arcsinh\sqrt{\varepsilon/(4a)}$, against $\sqrt{\varepsilon/a}=0.23408$.
\end{numcheck}
% src: L4_oz (a); C_v09 (e)

\begin{numcheck}[Ornstein--Zernike decay in two dimensions]\label{chk:sca-oz-lattice}
On a $512\times512$ torus, for the resolvent $(I-W)^{-1}$ at $c=0.2$: at $\varepsilon=0.05$ the ratio to the form \eqref{eq:sca-oz-prefactor} rose to $0.9963$, $0.9965$ and $0.9958$ along the axis, the diagonal and the direction $(2,1)$ at $r\approx30\xi$, and similarly at $\varepsilon=0.01$; the relative anisotropy of the support function was $0.00537$ against the small-$\varepsilon$ value $\varepsilon/(48a)=0.00556$ at $\varepsilon=0.05$, and $0.00105$ against $0.00105$ at $\varepsilon=0.01$; at $\varepsilon=0.3$ the support function gives $1.4254$ and $1.4806$ along the axis and the diagonal, and the local rates corrected for the prefactor approach $1.4260$ and $1.4822$. By Fourier inversion on a $1024\times1024$ torus at three parameter sets, the fitted pointwise axis rates were $0.51023$, $0.22421$ and $1.23420$ against $\arccosh(1+\varepsilon/(2a))=0.51083$, $0.22455$ and $1.23024$, and the diagonal rates $0.51260$, $0.22442$ and $1.26222$ against $\sqrt2\arccosh(1+\varepsilon/(4a))=0.51357$, $0.22478$ and $1.26655$. At $(c,a)=(0.1,0.2)$ in two dimensions the hyperplane-summed decay rate was $0.6931=\arccosh(1+\varepsilon/(2a))$, and \eqref{eq:sca-lengths-r3} with $\cosh m=1+\varepsilon/(2a)$ gave $\xi_2=1.3765$, the computed value, whereas the one-dimensional relation $\cosh m=(1-c)/(2a)$ would give $0.6156$.
\end{numcheck}
% src: L4_oz (b'); C_v09 (d); C2_b1_size (the D >= 2 repair)

\begin{numcheck}[Second-moment lengths]\label{chk:sca-lengths}
(a) By direct summation, $\xi_2=2.70327$ and $6.26909$ in one dimension and $1.91150$ and $4.43292$ in two at $(c,\varepsilon)=(0.2,0.05)$ and $(0.2,0.01)$, equal to \eqref{eq:sca-xi2-exact} to all printed digits; an independent summation gave $\xi_2^2=19.30303$ and $9.65152$ in one and two dimensions, equal to the formula, with $a/\varepsilon=19.5$ and $9.75$. (b) On a ring of 4096 sites with four random kernels of range one to four, the covariance-side and precision-side lengths agreed to all printed digits ($3.058728$, $23.469588$, $7.390093$, $2.808015$), the root $m_d$ of \cref{prop:sca-lengths-relation}(b) equaled the fitted decay rate to six digits, and \eqref{eq:sca-lengths-bracket} held, for example $23.4696\le23.4892\le23.6398$. For the medium of \cref{thm:sca-oz} in one dimension, $\xi_{\mathrm{OZ}}/\xi_2=1.041$, $1.003$ and $1.0004$ at $\xi_2=1.69$, $7.02$ and $15.79$. (c) In two dimensions the on-site variance grew with slope $0.19867$ in $\ln(1/\varepsilon)$ between $\varepsilon=10^{-3}$ and $10^{-4}$, against $q/(8\pi a)=0.19897$.
\end{numcheck}
% src: L4_oz (c); C_v09 (g); B1_size D; L4_time_scale (b)

\begin{numcheck}[Persistence]\label{chk:sca-persistence}
Three strongly non-normal $W$ with $\rho(W)=0.9$: $\Vert\tilde W\Vert$ equaled $(1-\lambda_{\min}(\Sigma^{-1/2}Q^\zeta\Sigma^{-1/2}))^{1/2}$ to five digits ($0.98319$, $0.98812$, $0.99996$), and $\Vert\tilde W^h\Vert\le\Vert\tilde W\Vert^h$ at $h=1$, $5$ and $20$, while $\Vert\tilde W^h\Vert$ tracked $\rho^h$ for large $h$; in the third case $\Vert\tilde W\Vert=0.99996$ while $\Vert\tilde W^{20}\Vert=0.47$.
\end{numcheck}
% src: L4_time_scale (d)

\begin{numcheck}[Autocovariance and diffusive scaling]\label{chk:sca-autocovariance}
In the medium of \cref{thm:sca-oz} with $c>0$, the local autocovariance was compared with the continuum approximation $\frac q2\int_h^\infty e^{-\varepsilon_\lambda u}(4\pi au)^{-D/2}\mathrm{d}u$, $\varepsilon_\lambda=\ln(1/\lambda)$: at $\varepsilon=0.002$ the ratio lay in $[0.994,1.005]$ for $h$ from 10 to 3000 in one dimension and in $[0.993,1.017]$ in two. That range has $\varepsilon h\le6$, where the integral form applies and the asymptote \eqref{eq:sca-autocov}, its limit for $\varepsilon h\gg1$ at small $\varepsilon$, is not yet accurate: at $\varepsilon=0.002$ and $h=3000$ the ratio of \eqref{eq:sca-autocov-def} to \eqref{eq:sca-autocov} was $0.932$ in one dimension and $0.872$ in two. The asymptote itself was tested at $c=0.2$ and $\varepsilon=0.05$, where the same ratio was $0.8715$, $0.9583$, $0.9886$ and $0.9969$ at $h=50$, $200$, $800$ and $3000$ in one dimension and $0.7703$, $0.9201$, $0.9775$ and $0.9938$ in two; the shortfalls at $h=3000$, $0.0031$ and $0.0062$, equal $D\lambda^2/((1-\lambda^2)h)$ to the digits shown, about $D/(2\varepsilon h)$. Local decay rates approached $\varepsilon_\lambda+D/(2h)$, for example $0.01021$ and $0.01037$ at $h=3000$ and $\varepsilon=0.01$ in one and two dimensions. The ratio $\xi_2^2/\tau$ was $0.39499$ against $a=0.395$ at $\varepsilon=0.01$ in one dimension, and $0.19750$ against $0.1975$ in two.
\end{numcheck}
% src: L4_time_scale (a)

\begin{numcheck}[The record model]\label{chk:sca-records}
(a) Over 3000 random reduced models with one to four records, the root of $g_z$ equaled $\lambda_{\min}(\mathsf R)$ to a relative error of $8.5\times10^{-13}$, and $\rho(I-\mathsf R)=1-\varepsilon_{\mathrm{rec}}$ to $2.9\times10^{-15}$ whenever $I-\mathsf R\ge0$; an independent computation found $1.3\times10^{-12}$ and $3.1\times10^{-15}$ in 2993 nonnegative cases. (b) No sign mismatch between $\varepsilon-\varepsilon_{\mathrm{rec}}$ and the criterion of \cref{prop:sca-records-memory}(c) occurred in 15962 stable one-record draws with $\upsilon$ zero, small or arbitrary in $[0,1)$, and no mismatch of \eqref{eq:sca-memory-criterion} in 20000 several-record draws; the brackets of \cref{prop:sca-records-memory}(a) and (b) held in 20000 draws. (c) Mixed records at $\varepsilon=0.002$ and $\upsilon=0$: the durable record $(\beta,\omega,\delta)=(0.05,0.01,10^{-4})$ alone gave $\varepsilon_{\mathrm{rec}}=0.000408$ and multiplied $\tau$ by $4.900$; with the perishable record $(0.2,0.1,0.01)$ added, $\varepsilon_{\mathrm{rec}}=0.002291$ and $\tau$ fell to $0.873$ of its value without records. (d) The exact reduction of \cref{prop:sca-records}(d) held to $4.1\times10^{-15}$ on random networks of 4 to 12 agents with one to three records and heterogeneous write weights. (e) Special cases: one record with $(\varepsilon,\beta,\omega,\delta)=(0.01,0.2,0.1,0)$ gave $\varepsilon_{\mathrm{rec}}=0.003260$ and $\tau=306.24$; in two turnover instances the secular roots equaled the exact identity rates $0.0012379$ and $0.0021905$. (f) For level-0 reading with fixed prior precision the displacement fraction equaled the sensing gain before reading in three instances ($0.166667$, $0.142857$, $0.019608$). (g) Thirty agents on a ring and one archive with retention $1-\delta$: $\rho(W)=0.9000$, $0.9862$ and $0.9951$ for $\delta=0.1$, $0.01$ and $0.001$, at least the archive's self-weights $0.894$, $0.984$ and $0.993$, and an agent's autocorrelation at lag 50 rose from $1.1\times10^{-4}$ to $5.2\times10^{-2}$.
\end{numcheck}
% src: B5_records parts A-E, G, H; C2_b5_records; L4_time_scale (d)

\begin{numcheck}[The length with records]\label{chk:sca-records-length}
A ring of 2000 sites with $a=0.2$ and $\varepsilon=0.002$, so that $\xi_2=9.995$ hops and $\tau=499.5$ rounds without records. For no record; one record $(\beta,\omega,\delta,\ell_r)=(0.1,0.01,10^{-4},0)$; $(0.1,0.1,5\times10^{-3},6)$; $(0.1,0.01,5\times10^{-3},6)$; $(0.1,0.01,10^{-4},6)$; and two records $(0.05,0.01,10^{-4},0)$ and $(0.05,0.1,5\times10^{-3},6)$, the real-space second moment gave $\xi_2=9.995$, $7.265$, $16.308$, $3.592$, $22.527$ and $12.957$ hops, equal to formula (a) of \cref{prop:sca-records-length} in every case. The slow-branch value $(2(1-x_0)\mathcal D_{\mathrm{eff}}/(x_0(2-x_0)))^{1/2}$, the leading term in the proof of (b), was $9.995$, $8.078$, $16.927$, $6.278$, $25.742$ and $14.333$, and the first-order form of \eqref{eq:sca-records-xi} $10.000$, $8.165$, $16.903$, $6.202$, $25.820$ and $14.832$; the correlation times were $499.5$, $3702.5$, $286.1$, $211.3$, $3702.5$ and $1381.1$. In every configuration \eqref{eq:sca-deff} agreed with the first-order perturbation of the slow eigenvalue to $4.2\times10^{-16}$. For the asymptotics of (b), the leaks of the second, fourth and sixth configurations, with $q_z=0.5$, were divided by $10$, $100$ and $1000$: the ratio of $\xi_2^2$ to $\mathcal D_{\mathrm{eff}}/\varepsilon_{\mathrm{rec}}$ tended to one, its relative error divided by $\varepsilon_{\mathrm{rec}}$ settling at $167.3$, $181.4$ and $123.2$, as the error term of (b) requires. The fitted decay lengths of the covariance in the second to fifth configurations were $8.18$, $17.06$, $6.86$ and $25.90$ hops. On the same ring, the agents' marginal world-conditional precision had no positive off-diagonal entry with no record, with a durable local record and no record noise, and with a perishable record read at range and no record noise; with the durable record and record noise of variance one it had eleven positive off-diagonal entries, the largest $7.6\times10^{-3}$.
\end{numcheck}
% src: B5_records part F; L5-P9 fitted lengths (via B5); C2_b5_records (S4); editor's Lyapunov computation for the asymptotics of (b) (double precision, one thread)

\begin{numcheck}[Hub component of a global record]\label{chk:sca-global-records}
One record read by all agents of a ring with weight $\beta$ and written uniformly with total weight $\omega$, no record noise, in the setting of \cref{rem:sca-global-records}. On a ring of 12 agents with $a=0.1$, $\varepsilon'=10^{-3}$, $\beta=0.2$, $\omega=0.05$ and $\delta=10^{-4}$, the agents' covariance from the discrete Lyapunov equation equaled the branch $p\neq0$ continued to all wavevectors plus the constant $[\hat\Sigma_{00}(0)-q/(1-(1-\varepsilon'-\beta)^2)]/N$ to within $2\times10^{-12}$. With $s=(\varepsilon'+\mathcal R\delta)(\beta+\omega)/(\beta\omega)$, the constant was the following fraction of $\hat\Sigma_{00}(0)/N$ at $(\varepsilon',\beta,\omega,\delta)$: $0.996$ at $(10^{-4},0.1,0.1,10^{-4})$, where $s=0.004$; $0.962$ at $(10^{-3},0.2,0.05,10^{-4})$, where $s=0.035$; $0.484$ at $(10^{-5},0.1,10^{-3},10^{-5})$, where $s=1.02$; $0.093$ at $(10^{-5},0.1,10^{-3},10^{-4})$ and $0.078$ at $(10^{-2},0.1,10^{-3},0)$, where $s=10.1$. In ten instances the fraction stayed within $0.015$ of $1/(1+s)$. The last instance has $\varepsilon_{\mathrm{rec}}\approx10^{-4}\ll\beta$, so smallness of the effective leak beside the read weight does not make the hub component dominant.
\end{numcheck}
% src: fix/08 scripts/final_global_record.py (scipy solve_discrete_lyapunov, double precision, one thread)

\begin{numcheck}[Algebraic decay]\label{chk:sca-longrange}
A ring of $2^{17}$ sites with $c=0.3$ and $\lambda=0.95$. The ratio of $\Sigma(x)$ to $2q\lambda a\mathsf J(x)/(1-\lambda^2)^2$ fell from $1.3034$ at $x=25$ to $1.0006$ at $x=3200$ for $\sigma=1.5$, and from $5.1427$ to $1.0000$ for $\sigma=3$, with fitted tail exponents $2.503$ and $4.000$ between $x=800$ and $3200$ (against $1+\sigma=2.5$ and $4$); for $\sigma=0.5$ it rose from $0.11$ to $0.76$ over the same range, with corrections decaying roughly as $x^{-0.48}$.
\end{numcheck}
% src: L4_longrange (a)

\begin{numcheck}[Fronts]\label{chk:sca-fronts}
The front $r_{0.1}(t)$ of $\Pr(|Y_t|\ge r)$ in one dimension, on rings of $2^{20}$ sites for $\sigma=0.5$ and $0.75$ and of $2^{17}$ sites for $\sigma=1$, $1.5$ and $3$, with the fronts at most $3.2$ percent of the half-length. Least-squares fits of $\ln r_{0.1}$ against $\ln t$ over doubling times, $t=1$ to $8$ for $\sigma=0.5$, $t=2$ to $32$ for $\sigma=0.75$ and $t=16$ to $512$ for the others, gave exponents $1.897$, $1.305$, $0.988$, $0.671$ and $0.469$ for $\sigma=0.5$, $0.75$, $1$, $1.5$ and $3$, against $2$, $1.333$, $1$, $0.667$ and $0.5$. The fits depend on the window: at $\sigma=3$, where the fronts are only $8$ to $39$ sites, the window $t=32$ to $256$ gave $0.497$, and at $\sigma=0.5$ the local slopes between successive times were $1.879$, $1.906$ and $1.901$, so the shortfall of five percent there did not close within the window. The single-big-jump ratio $\Pr(Y_t=x)/(t\Pr(X_1=x))$ at $\sigma=1.5$ and $t=16$ was $1.0133$, $1.0017$, $1.0002$ and $1.0000$ at $x=200$, $800$, $3200$ and $12800$.
\end{numcheck}
% src: L4_longrange (b), (c)

\begin{numcheck}[Gaussian broadcasting on trees]\label{chk:sca-trees}
The recursion of \cref{thm:sca-trees}(a) agreed with direct Gaussian conditioning to ten digits on binary trees of depth 7 (255 vertices), ternary trees of depth 5 (364 vertices) and Galton--Watson trees, with heterogeneous edge correlations and leaf noise $0$ or $0.3$. On $b$-ary trees: at $b\alpha^2=0.8$, $I_{40}=2.216\times10^{-5}$ against the asymptotic $2.215\times10^{-5}$; at $b\alpha^2=1$, $I_{40}=2.440\times10^{-2}$ against $2.500\times10^{-2}$; at $b\alpha^2=1.25$ and $1.8$ the limits $0.2554$ and $1.0986$. The closed form of the union bound in \cref{thm:sca-percolation} agreed with truncated sums at three parameter sets, for example $0.02000000$ at $\Delta s_e=0.1$, $s_m=0.5$, $d=3$.
\end{numcheck}
% src: L4_tree_info (a), (b), (d)

\begin{numcheck}[Gaussian information percolation]\label{chk:sca-gaussian-percolation}
Four hundred random trees with edge correlations in $[0.05,0.99]$, three leaf-noise levels and five source signal-to-noise ratios: the largest ratio of the left side of \eqref{eq:sca-gauss-perc} to the right side was $1.0000$, attained in the small-signal limit of a single edge, with no violation.
\end{numcheck}
% src: L4_tree_info (c)

\begin{numcheck}[Coherence volume and susceptibilities]\label{chk:sca-coherence-volume}
(a) Disks of radius 2, 4, 8 and 12 (13 to 441 agents) in the two-dimensional medium of \cref{thm:sca-oz} on a $128\times128$ torus, at $\varepsilon=0.05$ ($\xi_{\mathrm{OZ}}=1.96$, $\chi=7.0$) and $\varepsilon=0.01$ ($\xi_{\mathrm{OZ}}=4.45$, $\chi=28.4$): the best correlational coherence from the quadratic program equaled the smallest member correlation with $p^{\ast\top}X$ to four digits, and $|C|\theta^{\ast2}\le\lambda_{\max}(\mathrm{Corr}_C)\le\chi$ in every case; at $\varepsilon=0.01$ the disk of 441 agents has $\theta^\ast=0.19$. (b) For 180 random subsets of 60 random nonnegative networks the minimax value matched $\min_i\Corr(X_i,Z^\ast)$ to $1.5\times10^{-7}$, the chain of \cref{thm:sca-coherence-volume}(b) held, and the largest susceptibility was $0.19$ of the bound of \cref{cor:sca-size}(a). (c) On the complete graph at $\lambda=0.9$, $N\theta^{\ast2}=3.690$, $5.048$ and $5.241$ at $N=10$, $100$ and $1000$, against $1/(1-\lambda^2)=5.263$. (d) Lattice susceptibilities: $\chi/(2\xi_{\mathrm{OZ}})=0.81$, $0.88$ and $0.94$ in one dimension as $\varepsilon$ falls from $0.04$ to $0.0025$; in three dimensions ($c=0.2$, a $48^3$ torus) $\chi$ equaled $q\xi_2^2/(2a\lambda\Sigma(0))$ to rounding, as it must, since $\chi=q/((1-\lambda^2)\Sigma(0))$ and \cref{eq:sca-xi2-exact} make this an identity in every dimension; the form $q\xi_2^2/(2a\Sigma(0))$ of \cref{cor:sca-size}(d) is therefore $\lambda\chi$ and lies below $\chi$ by the relative amount $\varepsilon$ exactly, which is $4$, $1$, $0.6$ and $0.25$ percent at $\varepsilon=0.04$, $0.01$, $0.006$ and $0.0025$; in two dimensions the leading logarithmic form is far from the exact $\chi$ at these $\varepsilon$, so the exact $\Sigma(0)$ must be used. (e) A directed copy chain with weights $0.95$ has $\rho(W)=0$, so $\tau=0$, while $1/(1-\lambda_r\lambda_c)=10.26$ and $\chi=1$: the bound of \cref{cor:sca-size}(a) holds but is not a memory time.
\end{numcheck}
% src: L4_time_scale (b), (c); C_v10

\begin{numcheck}[Products of projections]\label{chk:sca-alternating}
(a) Two codes in six dimensions, with three-dimensional transmitted subspaces and a one-dimensional common subspace, in a common prior metric: $\Vert(P_{V_2}P_{V_1})^n-P_\cap\Vert=c^{2n-1}$ to seven digits for $n=1$ to $20$, with $c=0.683035$. (b) Three codes in cyclic order, four-dimensional subspaces with a common vector: $c_{\mathrm{cyc}}=0.963301$, and $\Vert A_{\mathrm{cyc}}^n-P_\cap\Vert\le c_{\mathrm{cyc}}^n$ at $n=1$, $5$, $10$, $20$ and $40$ (at $n=40$, $0.0755$ against $0.224$). (c) Quasi-periodic orders whose blocks are random permutations of the three codes: $c_\ast=0.96930$, and the bound of \cref{thm:sca-alternating}(b) held at every block.
\end{numcheck}
% src: L4_proj (1)-(3)

\begin{numcheck}[Lyapunov spectra of random relays]\label{chk:sca-lyapunov}
A six-dimensional world, rank-two codes $H=\diag(1,0.6)U^\top$ with $U$ a Haar-random $6\times2$ frame, message noise $0.1I$ and identity prior, so that the hop map is $A=UGU^\top$ with $G=\diag(1/1.1,0.36/0.46)$. The law of the hop map is invariant under conjugation by orthogonal matrices, so the uniform measure on $k$-planes is stationary and $\Gamma_1+\dots+\Gamma_k$ is the expectation of the logarithm of the $k$-volume of $AE_k$ for a fixed orthonormal $k$-frame $E_k$; the exponents are averages of this one-step quantity over $2\times10^6$ independent draws, and path products with repeated orthogonalization ($400$ paths of $300$ hops) agreed with them within their standard errors. Pure relays: exponents $-0.918$ and $-1.506$ per hop, gap $0.588$ with standard error $0.001$, in agreement with the values $\Gamma_1=-3/4+\ln((d_1+d_2)/2)=-0.9174$ and $\Gamma_1+\Gamma_2=\ln(d_1d_2)-25/12=-2.4238$ that the same one-step expectation gives in closed form for $G=\diag(d_1,d_2)$, whose gap is $0.5889$; abstraction length $1.70$ hops, so the ratio $\sigma_2/\sigma_1$ of a typical product falls below $0.01$ after $\ln100/0.588\approx8$ hops. The product of the $2\times2$ factors $GU_m^\top U_{m-1}$ of \cref{sec:sca-abstraction} gave the same two exponents within the standard errors, with sums $-2.424$ and $-2.427$, and the smallest $|\det U_m^\top U_{m-1}|$ among $1.2\times10^5$ draws was $5\times10^{-6}$, so every factor was invertible. The crossing hop scatters widely: over $2000$ independent realizations the second singular value was still above $0.01$ of the largest in $74$, $44$, $24$ and $4$ percent of them at $n=5$, $8$, $11$ and $20$, with median ratio $0.041$ at $n=5$ and $0.007$ at $n=8$. Receivers that keep part of their own belief, with hop map $(1-\omega)I+\omega A$ and $\omega=1/2$, so that the retention weight is $1/2$: six exponents from $-0.423$ to $-0.545$, gap $0.0296$ with standard error $0.0002$, abstraction length $33.8$ hops, and in one realization $6$, $5$ and $2$ singular values above the threshold at $n=10$, $20$ and $60$. A single path of $4000$ hops determines this gap only to about $\pm0.004$, hence the length only to within several hops. With identical frames and isotropic gains the rank stayed at the code rank, two.
\end{numcheck}
% src: L4_proj (6)
